\section{Why subgroups of free groups are free}\label{sec:schreier}
The overview gives this consequence as an example of topology proving a theorem of algebra. The topology here is a graph covering a bouquet of circles. Even when that graph is infinite, the part seen by a word can be computed within a bound set by the word itself.

\begin{theorem}[Schreier basis for a presented subgroup]\label{thm:schreier}\label{idx-schreier4}
Fix the free group $F_r$ of reduced words on $r$ letters and their inverses. For each fixed PR membership test $H(w)$ satisfying the subgroup and reduction laws, there is a PR test for basis symbols and mutually inverse PR procedures between elements of $H$ and reduced words on those symbols. In particular the subgroup is freely generated by a specified, possibly infinite, decidable set of words.
\end{theorem}
\begin{proof}
Use right cosets $Hw$ as notation for the numerical equivalence relation
$u\sim v$ if $uv^{-1}\in H$. Products always include free reduction. The subgroup laws give equivalence, and appending the same word preserves it. Among reduced $v\sim w$ with $|v|\le |w|$, choose the shortest, breaking ties lexicographically in a fixed order on signed letters. Call it $R(w)$. This is the least of the whole class because $w$ itself is a candidate. In particular $R$ is constant on classes. A scan of at most $(2r+1)^{|w|+1}$ candidates suffices (for $r=0$ only the empty word occurs).

Every prefix of a representative is again a representative. Otherwise replace that prefix by a smaller equivalent word and append the unchanged suffix. Free reduction can shorten it; if its length stays equal, the first lexical improvement persists. The result represents the same class and is smaller, a contradiction. The representatives therefore form a rooted tree under deleting the last letter, with root the empty word. Consider the graph with vertices these representatives and one oriented edge $(u,a)$ from $u$ to $v=R(ua)$ for each positive generator $a$. Traversing that edge backwards reads $a^{-1}$. The prefix tree consists exactly of edges for which $ua=v$ as reduced words: if cancellation occurs, the endpoint is the prefix of $u$; without cancellation, $u$ is the prefix of $v$. Each non-tree edge gives the basis symbol $(u,a)$ with value
\[
                          b(u,a)=uav^{-1}\in H .
\]
The test for a symbol is $R(u)=u$, $1\le a\le r$, and $ua\ne R(ua)$ as reduced words.

For a subgroup word $h$, start at the root and read it letter by letter. At a positive letter use the indicated edge; at $a^{-1}$ its positively oriented edge starts at $R(ua^{-1})$. Record nothing on tree edges and the signed symbol on every other edge; freely reduce this output. The final vertex is the root. A letter followed by its inverse traverses an edge and immediately returns on it, so removing this pair leaves the reduced output invariant. Two subgroup words concatenate as two loops at the root. Thus the reading map is a homomorphism. In the reverse direction substitute $b(u,a)$ for its symbol, using inverses when appropriate. The substituted word of a walk from $u_0$ to $u_k$ equals $u_0\,w\,u_k^{-1}$, since the factors at intermediate vertices cancel, and the omitted tree factors equal the identity. For a based loop this gives $h$ back. Conversely the path reading $b(u,a)$ follows the tree out along $u$, traverses the indicated non-tree edge and follows the tree home along $v^{-1}$. It records that one symbol. This proves mutual inverseness, generation and freedom. Distinct symbols have distinct group elements, since reading their expansions recovers the symbols.

At step $j$ of reading an input of length $L$ the representative has length at most $j$. The difference between testing a coset and deciding a real or infinitary predicate is important here: $H$ is a supplied numerical term. The proof of the reading identities is induction along a bounded word with tests on finite strings. It proves an iPRA schema with the stated membership and subgroup hypotheses.
\end{proof}

\begin{corollary}[The finite-index rank formula]\label{idx-schreier22}
Given a transitive right action of $F_r$ on $d\ge1$ numbered points by $r$ permutation tables, the stabilizer of a specified point is free of rank $1+d(r-1)$.
\end{corollary}
\begin{proof}
Send $Hw$ to the image of the basepoint under $w$. The test $uw^{-1}\in H$ says precisely that $u$ and $w$ send the basepoint to the same point. Transitivity gives every target and elementary path search gives a word of length at most $d-1$ reaching it. The graph has $d$ vertices and $dr$ positive edges, including loops and parallel edges. Its prefix spanning tree has $d-1$ edges; the remaining $dr-d+1$ give the displayed rank. When $r=0$, transitivity forces $d=1$ and the formula gives zero.
\end{proof}

\begin{theorem}[A finite list of subgroup generators suffices]\label{thm:fold}\label{idx-schreier29}
Given a finite list of words in $F_r$, one can compute a finite labelled graph deciding membership in their subgroup and a finite free basis. The construction is uniform and PR in the whole list.
\end{theorem}
\begin{proof}
Make a base vertex and a subdivided loop for each nonempty word, with its successive directed edges labelled by the input letters. An edge has a formal reverse carrying the inverse label. The reduced labels of closed based walks form exactly the generated subgroup. Indeed, cancelling immediate edge reversals leaves the label unchanged in $F_r$, and away from the base a walk without reversals must continue around the current loop until it returns. Conversely, concatenating the given loops and their reversals realizes each product.

If two distinct oriented edges of the same label start at the same vertex, identify them, their reverses, and their terminal vertices. Scan and repeat. This operation preserves the subgroup read by based walks. A walk downstairs lifts edge by edge to the old graph with possible jumps between the two merged terminal vertices. Bridge each jump by going backward along one old edge and forward along the other. This contributes a freely trivial label. The same bridge at the beginning or end aligns the lifted basepoint. If the terminal vertices already coincide, no bridges are needed. Conversely, projecting an old walk preserves its label. Each fold removes at least one unoriented edge, so their number is bounded by the sum $L$ of input lengths.

In the final graph each signed label has at most one outgoing edge per vertex. A based walk labelled by a word with adjacent inverse letters traverses an edge and its actual reverse at that pair: uniqueness holds at the intermediate vertex. Hence a reduced word is in the subgroup exactly when its unique attempted walk exists and ends at the base. This is a PR membership test. Choose a spanning tree; for each non-tree edge choose one orientation, take the tree path out, that edge, and the tree path home. These loops freely generate the edge-path group by the same cancellation and reading argument as above. The map from this group to $F_r$ is injective: a walk without immediate reversals has a freely reduced label, by determinism at each vertex. A nonempty reduced loop therefore cannot have trivial label. The labels of the constructed loops give a finite free basis of the prescribed subgroup. Breadth-first parent computation, each conflict search and relabelling, and the $L$ fold iterations are bounded finite operations.
\end{proof}
